% GENERATED FILE. Edit canonical lessons, then run npm run generate:books.
% canonical-source-hash: fnv1a64:d00e75c2fb283b44
% canonical-lessons: SA-C443-suryastah, SA-C443-astam-gacchati, SA-C443-godhulih, SA-C443-tamisra, SA-C443-naktam

\chapter{Sunset and Night}
\label{ch:sa-sunset-and-night}
\hlchaptermodality{443}

\begin{chapteropening}
\textbf{By the end of this chapter:} \emph{I can say when the sun sets and the dark night comes.}

Complete the last lesson of chapter 443: I can say when the sun sets and the dark night comes.
\end{chapteropening}

\section[\texorpdfstring{\sk{सूर्यास्तः} (sūryāstaḥ)}{सूर्यास्तः (sūryāstaḥ)}]{\sk{सूर्यास्तः} (sūryāstaḥ) — sunset}

\label{lesson:SA-C443-suryastah}



\begin{quote}
Before the new one: say the Sanskrit for the horizon, then the Sanskrit for dawns.
\end{quote}

\subsection*{You'll want to know: \sk{सूर्यास्तः}}
\textbf{\sk{सूर्यास्तः}} — \emph{sūryāstaḥ} — \textquotedblleft{}sunset\textquotedblright{}.

\begin{quote}
\sk{सूर्यास्ते} \sk{काकः} \sk{नीडं} \sk{गच्छति।}
\end{quote}

\begin{quote}
\emph{sūryāste kākaḥ nīḍaṃ gacchati.}
\end{quote}

\begin{quote}
At sunset the crow goes to its nest.
\end{quote}

A masculine noun.

\begin{cousinweb}
\textbf{\sk{सूर्यास्तः}} (\emph{sūryāstaḥ}) is built from \emph{sūrya}, \textquotedblleft{}the sun\textquotedblright{}, and \emph{asta}, \textquotedblleft{}home, setting\textquotedblright{}; MW gives the compound as sunset (MW).
\end{cousinweb}

\begin{grammarlens}[title={What you've built}]
One more word for getting around.
\end{grammarlens}

\subsection*{Your turn}
\begin{itemize}
\raggedright
  \item \emph{Say it:} \emph{sūryāstaḥ}
  \item \emph{Say it:} \emph{sūryāstaḥ}, once more
  \item \emph{Say it:} \emph{sūryāste kākaḥ nīḍaṃ gacchati.} — At sunset the crow goes to its nest.
  \item \emph{Recall:} say the Sanskrit for the horizon, then the Sanskrit for dawns, then say \emph{sūryāstaḥ} again
\end{itemize}

\begin{tcolorbox}[breakable,colback=teal!4,colframe=teal!35!black,title={Before you move on}]
What does \sk{सूर्यास्तः} mean? (\textquotedblleft{}Sunset\textquotedblright{}.) Say it once more.
\end{tcolorbox}
\section[\texorpdfstring{\sk{अस्तं} \sk{गच्छति} (astaṃ gacchati)}{अस्तं गच्छति (astaṃ gacchati)}]{\sk{अस्तं} \sk{गच्छति} (astaṃ gacchati) — sets in the west; goes down (of the sun)}

\label{lesson:SA-C443-astam-gacchati}



\begin{quote}
Before the new one: say the Sanskrit for dawns, then the Sanskrit for sunset.
\end{quote}

\subsection*{You'll want to know: \sk{अस्तं} \sk{गच्छति}}
\textbf{\sk{अस्तं} \sk{गच्छति}} — \emph{astaṃ gacchati} — \textquotedblleft{}sets in the west; goes down (of the sun)\textquotedblright{}.

\begin{quote}
\sk{सूर्यः} \sk{पश्चिमायाम्} \sk{अस्तं} \sk{गच्छति।}
\end{quote}

\begin{quote}
\emph{sūryaḥ paścimāyām astaṃ gacchati.}
\end{quote}

\begin{quote}
The sun sets in the west.
\end{quote}

Past (it …): \sk{सूर्यः} \sk{अस्तं} \sk{गतवान्}, \emph{sūryaḥ astaṃ gatavān}.

\begin{cousinweb}
\textbf{\sk{अस्तं} \sk{गच्छति}} (\emph{astaṃ gacchati}) is built from \emph{astam}, \textquotedblleft{}home, the setting of the sun\textquotedblright{}, and the root √gam, \textquotedblleft{}to go\textquotedblright{}: to go home; MW gives \emph{astam} with √gam as to set, of the sun; CDIAL 976 gives \emph{astam eti}, \textquotedblleft{}goes home, sets\textquotedblright{}, from the Rigveda on (MW; CDIAL 976).
\end{cousinweb}

\begin{grammarlens}[title={What you've built}]
One more word for getting around.
\end{grammarlens}

\subsection*{Your turn}
\begin{itemize}
\raggedright
  \item \emph{Say it:} \emph{astaṃ gacchati}
  \item \emph{Say it:} \emph{astaṃ gacchati}, once more
  \item \emph{Say it:} \emph{sūryaḥ paścimāyām astaṃ gacchati.} — The sun sets in the west.
  \item \emph{Recall:} say the Sanskrit for dawns, then the Sanskrit for sunset, then say \emph{astaṃ gacchati} again
\end{itemize}

\begin{tcolorbox}[breakable,colback=teal!4,colframe=teal!35!black,title={Before you move on}]
What does \sk{अस्तं} \sk{गच्छति} mean? (\textquotedblleft{}Sets in the west; goes down (of the sun)\textquotedblright{}.) Say it once more.
\end{tcolorbox}
\section[\texorpdfstring{\sk{गोधूलिः} (godhūliḥ)}{गोधूलिः (godhūliḥ)}]{\sk{गोधूलिः} (godhūliḥ) — dusk}

\label{lesson:SA-C443-godhulih}



\begin{quote}
Before the new one: say the Sanskrit for sunset, then the Sanskrit for sets in the west.
\end{quote}

\subsection*{You'll want to know: \sk{गोधूलिः}}
\textbf{\sk{गोधूलिः}} — \emph{godhūliḥ} — \textquotedblleft{}dusk\textquotedblright{}.

\begin{quote}
\sk{गोधूलौ} \sk{गोपालः} \sk{धेनुं} \sk{गृहम्} \sk{आनयति।}
\end{quote}

\begin{quote}
\emph{godhūlau gopālaḥ dhenuṃ gṛham ānayati.}
\end{quote}

\begin{quote}
At dusk the cowherd brings the cow home.
\end{quote}

A feminine noun.

\begin{cousinweb}
\textbf{\sk{गोधूलिः}} (\emph{godhūliḥ}) is built from \emph{go}, \textquotedblleft{}a cow, the earth\textquotedblright{}, and \emph{dhūli}, \textquotedblleft{}dust\textquotedblright{}: MW glosses it \textquotedblleft{}earth-dust\textquotedblright{}, the time of day when a haze seems to rise from the ground, and gives it as a time of twilight (MW).
\end{cousinweb}

\begin{grammarlens}[title={What you've built}]
One more word for getting around.
\end{grammarlens}

\subsection*{Your turn}
\begin{itemize}
\raggedright
  \item \emph{Say it:} \emph{godhūliḥ}
  \item \emph{Say it:} \emph{godhūliḥ}, once more
  \item \emph{Say it:} \emph{godhūlau gopālaḥ dhenuṃ gṛham ānayati.} — At dusk the cowherd brings the cow home.
  \item \emph{Recall:} say the Sanskrit for sunset, then the Sanskrit for sets in the west, then say \emph{godhūliḥ} again
\end{itemize}

\begin{tcolorbox}[breakable,colback=teal!4,colframe=teal!35!black,title={Before you move on}]
What does \sk{गोधूलिः} mean? (\textquotedblleft{}Dusk\textquotedblright{}.) Say it once more.
\end{tcolorbox}
\section[\texorpdfstring{\sk{तमिस्रा} (tamisrā)}{तमिस्रा (tamisrā)}]{\sk{तमिस्रा} (tamisrā) — a dark night}

\label{lesson:SA-C443-tamisra}



\begin{quote}
Before the new one: say the Sanskrit for sets in the west, then the Sanskrit for dusk.
\end{quote}

\subsection*{You'll want to know: \sk{तमिस्रा}}
\textbf{\sk{तमिस्रा}} — \emph{tamisrā} — \textquotedblleft{}a dark night\textquotedblright{}.

\begin{quote}
\sk{तमिस्रायां} \sk{दीपं} \sk{विना} \sk{गन्तुं} \sk{कठिनम्।}
\end{quote}

\begin{quote}
\emph{tamisrāyāṃ dīpaṃ vinā gantuṃ kaṭhinam.}
\end{quote}

\begin{quote}
On a dark night it is hard to go without a lamp.
\end{quote}

A feminine noun.

\begin{cousinweb}
\textbf{\sk{तमिस्रा}} (\emph{tamisrā}) is built on \emph{tamas}, \textquotedblleft{}darkness\textquotedblright{}, with an ending Pāṇini (5.2.114) gives; MW gives \emph{tamisrā} as a dark night (MW).
\end{cousinweb}

\begin{grammarlens}[title={What you've built}]
One more word for getting around.
\end{grammarlens}

\subsection*{Your turn}
\begin{itemize}
\raggedright
  \item \emph{Say it:} \emph{tamisrā}
  \item \emph{Say it:} \emph{tamisrā}, once more
  \item \emph{Say it:} \emph{tamisrāyāṃ dīpaṃ vinā gantuṃ kaṭhinam.} — On a dark night it is hard to go without a lamp.
  \item \emph{Recall:} say the Sanskrit for sets in the west, then the Sanskrit for dusk, then say \emph{tamisrā} again
\end{itemize}

\begin{tcolorbox}[breakable,colback=teal!4,colframe=teal!35!black,title={Before you move on}]
What does \sk{तमिस्रा} mean? (\textquotedblleft{}A dark night\textquotedblright{}.) Say it once more.
\end{tcolorbox}
\section[\texorpdfstring{\sk{नक्तम्} (naktam)}{नक्तम् (naktam)}]{\sk{नक्तम्} (naktam) — by night}

\label{lesson:SA-C443-naktam}



\begin{quote}
Before the new one: say the Sanskrit for dusk, then the Sanskrit for a dark night.
\end{quote}

\subsection*{You'll want to know: \sk{नक्तम्}}
\textbf{\sk{नक्तम्}} — \emph{naktam} — \textquotedblleft{}by night\textquotedblright{}.

\begin{quote}
\sk{उलूकः} \sk{नक्तं} \sk{चरति।}
\end{quote}

\begin{quote}
\emph{ulūkaḥ naktaṃ carati.}
\end{quote}

\begin{quote}
The owl moves about by night.
\end{quote}

\begin{cousinweb}
\textbf{\sk{नक्तम्}} (\emph{naktam}) is the word MW gives for by night, often set against \emph{divā}, \textquotedblleft{}by day\textquotedblright{}; MW gives no derivation, so its origin is uncertain (MW).
\end{cousinweb}

\begin{grammarlens}[title={What you've built}]
One more word for getting around.
\end{grammarlens}

\subsection*{Your turn}
\begin{itemize}
\raggedright
  \item \emph{Say it:} \emph{naktam}
  \item \emph{Say it:} \emph{naktam}, once more
  \item \emph{Say it:} \emph{ulūkaḥ naktaṃ carati.} — The owl moves about by night.
  \item \emph{Recall:} say the Sanskrit for dusk, then the Sanskrit for a dark night, then say \emph{naktam} again
\end{itemize}

\begin{tcolorbox}[breakable,colback=teal!4,colframe=teal!35!black,title={Before you move on}]
What does \sk{नक्तम्} mean? (\textquotedblleft{}By night\textquotedblright{}.) Say it once more.
\end{tcolorbox}
